\section{Economic interpretation and the continuing NBO contribution}\label{sec:discussion126}

The decision output of NBO is a feasible investment with its continuation
controller under a specified recursive preference. The relevant loss is the
value of that complete policy relative to the original optimum. A useful
learned representation must reduce this loss, or reduce the complete resources
required to reach a common economic tolerance. Those requirements motivate the
three current methodological results: action-aligned stopping, verified use of
the complete paid support record, and reuse of exact certificate geometry.
They retain the original Bellman problem rather than optimizing a different
fitting criterion and calling the result an economic improvement.

The combined evidence shows why all three accounts are necessary.
Production certificates meet the declared stronger regret target throughout the
522-case direct-service census. The independent comparison establishes a lower
mean original policy cost relative to
the completed fused control by a certified amount between $0.0003530$ and
$0.0003681$ on this fixed equally weighted census. There are nevertheless 13
individually certified losses against that control, and the stronger warm
adaptive comparator has an unresolved mean comparison, with 24 direct gains
and 26 direct losses. These are objective units of the stated investment model,
not a measured currency benefit or a population treatment effect. Neither a
lower upper bound on regret nor an internal accepted gain determines the
ranking of two completed competing policies.

The uncached computation further shows that economic information can be
purchased cheaply but processed expensively. Its substantial reduction in
terminal integration rows comes with a larger measured service-time bill than
warm adaptation. Exact support geometry was repeatedly rebuilt for questions
that shared the same paid information. The separately prospective cached
algorithm is motivated by that observed mechanism, retains the first result,
and gives the same optimization to every conventional and neural service.
Its complete paired costs and full scientific equality checks determine the
scope of its numerical improvement. Classical numerical implementation remains
a substantive comparator even when the actor is represented by a neural network.

The full-support solver answers a different question. It can exploit constraints
that a cheaper derivative or endpoint calculation leaves unused, and its exact
cut certificate retains validity even at an unsuccessful numerical exit or a
finite iteration cap. That additional information tightens 48 executed contrast
bounds and six candidate regret bounds. The tested candidate policies still
fail to acquire any new qualified strict gain. The extra LP and rational
arithmetic therefore have a measured bill and no additional realized policy
benefit in that study. This is a useful boundary for algorithm design: solving
a more complete information problem is economically warranted only when its
remaining decision value can justify its full cost. The conditional stopping
proposition makes that requirement precise without inventing a calibrated
price of computation.

Independent fitting addresses another part of the method. It reconstructs
finite-design objectives from the acquired labels, solves them through
separately assembled numerical routes, and verifies rational objective and
rounding bounds. Nonunique robust fits can have different coefficients without
contradicting the same finite-design objective. None of these facts makes a
finite training design into a uniform approximation theorem or supplies new
population policy evidence. The original actors are retained, and the shared
HiGHS library, unchanged labels and fixed features delimit the independence
actually achieved.

Exact ReLU and native hinge equivalence is equally consequential. With the same
features, readout, labels, feasible decisions and certification service, they
implement the same economic rule. Their shared geometry cannot establish a
neural-specific advantage. The original 96 sufficient neural adoption regions
are empty and have identical-policy conventional cost blockers under the
recorded full cost ranges. The follow-up applies its implementation change to
both sides of that comparison and reports its own adoption conditions. A
competitive neural contribution must survive these comparators, as well as
appropriate analytic, semianalytic, spline, polynomial and dynamic-programming
methods with their actual information and acquisition routes charged.

The mathematical scope is broader than this one computation but remains
conditional on explicit premises. The own-future construction and recursive
error accounts apply across their stated finite-horizon economic domains.
The stochastic-control and nonlinear-recursive-utility bridges require their
own drivers, filtrations, existence and boundary conditions. The complete
support program and its decision-resolution bounds concern scalar, smooth,
strongly convex original action objectives under a specified information
interface. They do not exclude algorithms using additional primitive formulas,
prove a randomized multivariate information lower bound, or establish neural
representation superiority at a common unrestricted resource budget.

The numerical application has two decision dates, two state coordinates,
one constrained control and one continuous innovation factor. Its state mesh,
risk levels and primitives are fixed; it is not newly calibrated to a firm or
a market. A substantive next extension must demonstrate action-relevant
information acquisition and complete decision value as these dimensions,
horizon and constraint interactions grow. A new learned actor's population
claim requires its own independent sample, and a practical adoption claim
requires an explicit decision workload and defensible resource value. The
present algorithms, proofs, reproducible positive and negative results, and
protected historical protocols provide the basis for those further tests
within the Neural Bellman Operators programme and its Econometrica objective.
